\section{A linear bound for the endpoint two}
\label{sec:endpoint-reduction}

The two possible low integer roots have different arithmetic constraints.
The endpoint two is confined to a linear region in the minimum exponent;
this gives a new nonintegral tail for even exponent triples.

\begin{theorem}
\label{thm:endpoint-two-bound}
Let $4\leq a\leq b\leq c$, and put
\[
R(a)=\frac{7a^2-2a+8}{a+2}=7a-16+\frac{40}{a+2}.
\]
If $h_C(2)=0$, then $c<R(a)$. In fact, $c\geq R(a)$ implies
$h_C(2)>0$.
\end{theorem}

\begin{proof}
Fix $a,c$ and expand the symmetric endpoint polynomial in $b$:
\[
h_C(2)=A(a,c)b^3+B(a,c)b^2+D(a,c)b+E(a,c),
\]
where
\begin{align*}
A(a,c)&=ac(a+c-1)+2a(a-1)+2c(c-1)-4,\\
B(a,c)&=c^2\bigl((a+2)c-(7a^2-2a+8)\bigr)\\
&\quad +(a^3+2a^2-6a-4)c+2(a-2)(a^2-2a-6),\\
D(a,c)&=(a-2)(c-2)\bigl(a^2c+a^2+ac^2+6ac+4a+c^2+4c-12\bigr),\\
E(a,c)&=2(a-2)(c-2)(a+c-2)(ac+a+c-2).
\end{align*}
For $a,c\geq4$, the coefficients $A,D,E$ are strictly positive.
Both $a^3+2a^2-6a-4$ and $a^2-2a-6$ are positive at $a=4$
and increasing on $[4,\infty)$. Thus $c\geq R(a)$ also gives $B>0$,
which proves the strict endpoint sign. Conversely, an endpoint zero
requires $B<0$, so $(a+2)c<7a^2-2a+8$.
\end{proof}

\begin{corollary}
\label{cor:even-linear-tail}
If $4\leq a\leq b\leq c$ are all even and $c\geq R(a)$,
then $G_{p^a q^b r^c}$ is not Laplacian integral.
For minimum exponent at least eight, $c\geq7a-12$ suffices.
\end{corollary}

\begin{proof}
Every entry of $C$ is even, so every integer eigenvalue of $C$ is even
by the rational-root theorem applied to $C/2$. In particular, one is
not an eigenvalue. Theorem~\ref{thm:endpoint-two-bound} excludes two,
and Theorem~\ref{thm:uniform-low-root} provides a noninteger root in
$(0,3)$. Equivalently, $h_C(0)<0<h_C(2)$ directly yields a root in
$(0,2)$, which cannot equal one. Its graph lift is noninteger.
For $a\geq8$, $40/(a+2)\leq4$ gives $R(a)\leq7a-12$.
\end{proof}

The triples $(8,10,44)$, $(10,14,58)$ and $(40,42,266)$ meet this
criterion with gcd two. Their largest exponents are below the earlier
quadratic cutoff and their spans are outside the earlier balanced region.
The already settled repeated triple $(10,10,12)$ has $h_C(2)=0$ and
obeys the strict bound; an integer endpoint does not establish integrality.

\begin{proposition}
\label{prop:endpoint-divisors}
For fixed $a,b\geq4$, any positive integer $c$ on an endpoint-zero
surface satisfies
\begin{align*}
h_C(1)=0&\ \Longrightarrow\ c\mid(a-1)(b-1)(a+b-1)(ab+a+b-1),\\
h_C(2)=0&\ \Longrightarrow\ c\mid2(a-2)(b-2)(a+b-2)(ab+a+b-2).
\end{align*}
\end{proposition}

\begin{proof}
Both endpoints are integer cubics in $c$. Their constant terms are
the displayed nonzero products. Reduce each endpoint equation modulo $c$.
\end{proof}

For a fixed pair this gives finite divisor candidates, which must still
satisfy the endpoint equation. Reza's requested diagnostic for middle
exponents through fifteen now has a complete independent certificate.

\begin{corollary}
\label{cor:middle-fifteen}
Every positive ordered three-prime exponent triple $a\leq b\leq c$
with $b\leq15$ gives a nonintegral ideal intersection graph.
This conclusion uses the complete finite endpoint certificate below.
\end{corollary}

\begin{proof}
Repeated exponents are covered by Theorem~\ref{thm:repeated}, and
minimum exponents one, two and three by Theorems~\ref{thm:unit},
\ref{thm:minimum-two} and~\ref{thm:minimum-three}, respectively.
It remains to consider $4\leq a<b\leq15$ and $c>b$.

For each of these $66$ pairs, both endpoint polynomials are integer
cubics in $c$ with the positive constant terms in
Proposition~\ref{prop:endpoint-divisors}. The
\href{https://github.com/the-omega-institute/ideal-intersection-laplacian/blob/10ff01f0625c7264b8bc5bb33b1a9d7dcb88a93c/results/endpoint-surfaces-verification.json}{complete endpoint certificate}
records all $132$
coefficient lists and exhausts their $8\,527$ positive divisors using
integer trial division. Independent integer Horner evaluation at every
one of the $7\,585$ divisors above $b$ gives no zero. Thus neither
endpoint vanishes for any integer $c>b$, with no imposed upper bound
on $c$. Theorem~\ref{thm:uniform-low-root} supplies a positive root
in $(0,3)$; it can equal neither one nor two, proving nonintegrality.
\end{proof}

The reflected quotient in this diagnostic uses
$N=|V(H)|=a+b+c+ab+ac+bc$, not the full graph vertex count.
The inclusive requested bound $b\leq15$ and exact rational linear-factor
extraction are cross-checked against the independent divisor certificate.
The empty lists settle this pair range; they do not prove either
endpoint surface globally empty.

The remaining all-even triples have gcd two,
$8\leq a<b<c<R(a)$ and $b\geq16$, outside the preceding proved regions.
The mixed-parity surface $h_C(1)=0$ still needs another obstruction.
